\documentclass[11pt]{article}

% --- Preamble: the kind of setup you'd reuse across posts ---
\usepackage{amsmath, amssymb, amsthm}

\newtheorem{theorem}{Theorem}
\newtheorem{lemma}{Lemma}
\theoremstyle{remark}
\newtheorem*{remark}{Remark}

% Custom macros -- the real test of a converter
\newcommand{\ev}[1]{\left\langle #1 \right\rangle}
\newcommand{\dd}{\mathrm{d}}
\newcommand{\R}{\mathbb{R}}
\newcommand{\grad}{\nabla}
\newcommand{\Ham}{\mathcal{H}}

\title{The Virial Theorem Is a Statement About Scaling}
\author{Jan Glowacki}
\date{July 2026}

\begin{document}
\maketitle

\begin{abstract}
The virial theorem is usually introduced as a time-averaging trick and then
filed away. I want to argue that its natural home is the theory of scaling
symmetries, where it becomes a one-line consequence of stationarity under
dilations. Nothing here is new; I just think the scaling viewpoint deserves
more airtime than it gets.
\end{abstract}

\section{The usual story}

For a bounded system of particles with positions $\mathbf{r}_i$ and momenta
$\mathbf{p}_i$, one defines the virial $G = \sum_i \mathbf{p}_i \cdot \mathbf{r}_i$
and observes that its long-time average $\dot{G}$ vanishes. Writing
$T$ for the kinetic energy, this yields the familiar
\begin{equation}
  2\,\ev{T} = -\, \Big\langle \sum_i \mathbf{F}_i \cdot \mathbf{r}_i \Big\rangle ,
  \label{eq:virial}
\end{equation}
and for a potential $V \propto r^{n}$ one gets the clean relation
$2\ev{T} = n\,\ev{V}$. This is correct, useful, and — I think — slightly
misleading about \emph{why} it is true.

\section{The scaling viewpoint}

Consider instead the response of the action to a uniform rescaling of
coordinates, $\mathbf{r}_i \mapsto \lambda\, \mathbf{r}_i$. Let the potential be
homogeneous of degree $n$, so that $V(\lambda \mathbf{r}) = \lambda^{n} V(\mathbf{r})$.
Then the kinetic and potential pieces of the Hamiltonian $\Ham = T + V$
scale oppositely, and stationarity of the on-shell action under $\lambda$
at $\lambda = 1$ is exactly Eq.~\eqref{eq:virial}.

\begin{theorem}[Virial from dilation invariance]
Let $V:\R^{3N}\to\R$ be homogeneous of degree $n$, and let the motion be
bounded. Then the time-averaged energies satisfy
\begin{equation}
  2\,\ev{T} = n\,\ev{V}.
  \label{eq:scaling}
\end{equation}
\end{theorem}

\begin{proof}
Euler's identity for a degree-$n$ homogeneous function gives
$\sum_i \mathbf{r}_i \cdot \grad_i V = n\,V$. Since $\mathbf{F}_i = -\grad_i V$,
the right-hand side of Eq.~\eqref{eq:virial} equals $n\,\ev{V}$. The chain of
equalities is short:
\begin{align}
  2\,\ev{T}
    &= -\, \Big\langle \sum_i \mathbf{F}_i \cdot \mathbf{r}_i \Big\rangle \\
    &= \Big\langle \sum_i \mathbf{r}_i \cdot \grad_i V \Big\rangle \\
    &= n\,\ev{V},
\end{align}
where the last step is Euler's identity under the time average.
\end{proof}

\begin{remark}
The point I care about: Eq.~\eqref{eq:scaling} did not require solving any
equations of motion. It is a \emph{Ward identity} for the (broken) dilation
symmetry, and the ``time average'' is doing the same job that the vacuum
expectation value does in the field-theory version.
\end{remark}

\section{Why I think this deserves more airtime}

Once you see the virial theorem as a scaling identity, three things that look
like separate facts collapse into one. The Kepler relation $2\ev{T} = -\ev{V}$
(here $n=-1$), the equipartition-flavored result for the harmonic trap
($n = 2$, giving $\ev{T} = \ev{V}$), and the tracelessness conditions that show
up in conformal field theory are all the same statement evaluated at different
homogeneity degrees. That unification is, to me, worth more than the trick.
I take up the conformal-field-theory side of this in \postlink{scaling-and-cft}{a companion note}.

\end{document}
